\begin{helper}
Fix one support-label pair \(B,J\), after the fixed-\((B,J)\)
blue-first/red-second residual scan has already supplied released residual
blocks.  Let \(U\) be the finite terminal coordinate set, let
\(\mathcal P\) be the finite transcript family, and let
\(\mathsf{Realized}(\beta,\tau)\) be the realized branch/transcript
predicate.  For a transcript \(\tau\), write \(P_\tau,A_\tau,Z_\tau\) for
its forced-present, forced-absent, and selected-absence coordinate sets.

Assume there is a nonnegative residual assignment mass \(\mu(A)\) on
terminal assignments \(A\subseteq U\), with total mass at most \(1\):
\[
\sum_{A\subseteq U}\mu(A)\le 1.
\]
For each realized \((\beta,\tau)\) with \(\tau\in\mathcal P\), assume a
released residual block
\[
\mathcal R_{\beta,\tau}\subseteq \{A:A\subseteq U\}
\]
and a block mass
\[
\rho_{\beta,\tau}
  =\sum_{A\in\mathcal R_{\beta,\tau}}\mu(A).
\]
Assume distinct realized pairs have disjoint released blocks.  Finally,
assume the fixed-\((B,J)\) residual cylinder domination already proved by
the released-block support:
\[
p^{|P_\tau|-u_R-u_B}(1-p)^{|A_\tau|-|Z_\tau|}
  \le \rho_{\beta,\tau}
\]
for every realized allowed pair.

Then there is a finite residual leaf set \(\Lambda_{B,J}\), a residual leaf
map \(\iota(\beta,\tau)\in\Lambda_{B,J}\), and nonnegative leaf masses
\(\nu:\Lambda_{B,J}\to\mathbb R\) such that
\[
\sum_{\lambda\in\Lambda_{B,J}}\nu(\lambda)\le1,
\]
the leaf map is injective on realized allowed pairs, and every realized
allowed pair satisfies
\[
p^{|P_\tau|-u_R-u_B}(1-p)^{|A_\tau|-|Z_\tau|}
  \le \nu(\iota(\beta,\tau)).
\]
\end{helper}
\begin{proof}
Apply the released-cylinder packing bound
\(\noderef{FixedSetHistoryCellRedReleasedCylinderPacking}\) to the complete
residual assignment space \(2^U\), the mass \(\mu\), and the released
blocks \(\mathcal R_{\beta,\tau}\).  The hypotheses match exactly: \(\mu\)
is nonnegative on \(2^U\), its total mass is at most \(1\), every realized
block is a subset of \(2^U\), \(\rho_{\beta,\tau}\) is the \(\mu\)-sum over
that block, and distinct realized blocks are disjoint.  Therefore
\[
\sum_\beta
  \sum_{\tau\in\mathcal P:\mathsf{Realized}(\beta,\tau)}
    \rho_{\beta,\tau}
\le 1 .
\]

Take the residual leaf set to be the finite product of branch labels and
transcript labels.  Define
\[
\iota(\beta,\tau)=(\beta,\tau).
\]
Give the leaf \((\beta,\tau)\) mass \(\rho_{\beta,\tau}\) when
\(\tau\in\mathcal P\) and \((\beta,\tau)\) is realized, and give it mass
\(0\) otherwise.  Nonnegativity follows from the block-mass identity and the
nonnegativity of \(\mu\) on each released assignment.  The total leaf mass is
the displayed packed sum over realized allowed pairs, hence is at most
\(1\).

If two realized allowed pairs have the same leaf, then the ordered pairs
\((\beta,\tau)\) are equal, so both the branch labels and transcripts are
equal.  Finally, for a realized allowed pair, the mass of its leaf is
\(\rho_{\beta,\tau}\), and the assumed residual cylinder domination is
exactly
\[
p^{|P_\tau|-u_R-u_B}(1-p)^{|A_\tau|-|Z_\tau|}
  \le \nu(\iota(\beta,\tau)).
\]
This constructs the fixed-\((B,J)\) residual leaf certificate from the
released-block data, rather than from realized-cell geometry alone.
\end{proof}
